% ============================================================================
%  SECTION 10 — CONCLUSION
%  >>> A REMPLIR PAR LES MEMBRES 1 & 4 <<<
%
%  Contenu attendu :
%    - Bilan : on a etudie une technique de resolution de programmes convexes
%      non lisses (gradient proximal), prouve ses garanties (O(1/k), O(1/k^2)),
%      compare aux alternatives (section 7), illustre sur l'audio (section 8).
%    - Rappeler le positionnement : FISTA = coeur, audio = application.
%    - Ouvertures : FISTA a restart, pas adaptatif (backtracking), extension
%      aux problemes non convexes.
% ============================================================================
\section{Conclusion}
\label{sec:conclusion}

% --- [MEMBRES 1 & 4] Ecrire la conclusion ici. ---
% TODO
